% Part: lambda-calculus
% Chapter: lambda-definability
% Section: introduction

\documentclass[../../../include/open-logic-section]{subfiles}

\begin{document}

\olfileid{lam}{ldf}{int}
\olsection{مقدمه}

در نگاه نخست، حساب لامبدا صرفاً حسابی بسیار انتزاعی از عبارت‌هایی است
که توابع و اعمال آن‌ها بر یکدیگر را نمایش می‌دهند. هیچ‌چیز در نحو حساب
لامبدا نشان نمی‌دهد که این‌ها توابعی روی نوع خاصی از اشیا هستند؛ به‌ویژه،
در نحو آن هیچ اشاره‌ای به اعداد طبیعی نشده است. عملیات پایهٔ آن، یعنی
اعمال و انتزاع‌های لامبدا، بر هر تابعی اعمال می‌شوند، نه فقط بر توابع
اعداد طبیعی.

بااین‌همه، با اندکی ابتکار می‌توان توابع حسابی، یعنی توابع روی اعداد
طبیعی، را در حساب لامبدا تعریف کرد. برای این کار، به‌ازای هر عدد طبیعی
$n \in \Nat$، ترم~$\lambd$ ویژه‌ای چون~$\num n$، یعنی \emph{عددنمای
چرچ}~$n$، تعریف می‌کنیم. (عددنماهای چرچ به افتخار آلونزو چرچ چنین
نامیده شده‌اند.)

\begin{defn}
  اگر~$n \in \Nat$، \emph{عددنمای چرچ} متناظر~$\num{n}$، عدد~$n$ را
  نمایش می‌دهد:
  \[
    \num{n} \ident \lambd[fx][f^n(x)]
  \]
  در اینجا~$f^n(x)$ حاصل اعمال~$f$ بر~$x$ به تعداد~$n$ بار است. برای
  مثال، $\num{0}$ برابر~$\lambd[fx][x]$ و~$\num{3}$ برابر
  $\lambd[fx][f(f(f\,x))]$ است.
\end{defn}
  
عددنمای چرچ~$\num n$ به‌صورت ترمی لامبدا کدگذاری می‌شود که تابعی با دو
آرگومان~$f$ و~$x$ را نمایش می‌دهد و~$f^n(x)$ را بازمی‌گرداند.
عددنماهای چرچ آشکارا در صورت نرمال‌اند.

نمایشی از اعداد طبیعی در حساب لامبدا البته تنها هنگامی سودمند است که
بتوانیم با آن‌ها محاسبه کنیم. محاسبه با عددنماهای چرچ در حساب لامبدا
یعنی اعمال ترم~$\lambd$ای چون~$F$ بر یک عددنمای چرچ، و کاهش ترم مرکب
$F\, \num n$ به صورت نرمال. اگر این ترم همواره به صورت نرمال کاهش یابد
و صورت نرمال همواره عددنمای چرچی چون~$\num m$ باشد، می‌توانیم خروجی
محاسبه را عدد~$m$ بدانیم. در آن صورت می‌توان~$F$ را معرف تابعی چون
$f\colon \Nat \to \Nat$ دانست؛ همان تابعی که~$f(n) = m$ اگر و تنها اگر
$F\, \num n \red \num m$. بنا بر خاصیت چرچ--راسر، صورت‌های نرمال، اگر
وجود داشته باشند، یکتا هستند. پس اگر~$F\, \num n \red \num m$، هیچ
ترم دیگری در صورت نرمال، و به‌ویژه هیچ عددنمای چرچ دیگری، وجود ندارد که
$F \, \num n$ به آن کاهش یابد.

برعکس، با داشتن تابعی چون~$f\colon \Nat \to \Nat$، می‌توانیم بپرسیم آیا
ترمی چون~$F$ وجود دارد که~$f$ را به این شیوه تعریف کند. در این صورت
می‌گوییم~$F$، تابع~$f$ را~!!{lambda define}s و~$f$ نیز
!!{lambda definable} است. این مفهوم را می‌توان به توابع چندموضعی و جزئی
تعمیم داد.

\begin{defn}
فرض کنید~$f\colon \Nat^k \rightharpoonup \Nat$ یک تابع جزئیِ
$k$-موضعی باشد. می‌گوییم ترم لامبدای~$F$، تابع~$f$ را
\emph{!!{lambda define}s} اگر برای همهٔ~$n_0$، \dots،~$n_{k-1}$،
\[
F \, \num{n_0} \, \num{n_1} \dots \num{n_{k-1}} \red \num{f(n_0, n_1, \dots,
  n_{k-1})}
\]
هرگاه~$f(n_0, \dots, n_{k-1})$ تعریف‌شده باشد؛ و در غیر این صورت
$F \, \num{n_0} \, \num{n_1} \dots \num{n_{k-1}}$ هیچ صورت نرمالی
نداشته باشد.
\end{defn}

نمونه‌ای بسیار ساده، توابع ثابت‌اند. ترم
$C_k \ident \lambd[x][\num{k}]$ تابع~$c_k\colon \Nat \to \Nat$ را
!!{lambda define}s، به‌طوری‌که~$c_k(n) = k$. زیرا برای هر~$n$،
$C_k \, \num n \ident (\lambd[x][\num{k}])\num n \redone \num{k}$.
تابع همانی به‌وسیلهٔ~$\lambd[x][x]$ !!{lambda defined} است. تعریف توابع
پیچیده‌تر البته دشوارتر است و اغلب ابتکار فراوان می‌طلبد. ازاین‌رو شاید
شگفت‌آور باشد که هر تابع محاسبه‌پذیر~!!{lambda definable} است. عکس این
ادعا نیز درست است: اگر تابعی~!!{lambda definable} باشد، محاسبه‌پذیر است.

\end{document}
